\documentclass{article}
\title{A Synthetic Document With Real Cross-References}
\author{Test Fixture}

\begin{document}
\maketitle

\begin{abstract}
This fixture exists to exercise the two code paths the five corpus documents
never reach: cross-references between blocks, and the merge rule for very
small blocks.\end{abstract}

\section{Background}
\label{sec:background}

Autoregressive models factor a joint distribution over a token sequence into
one conditional distribution per position. Each factor is evaluated from a
fixed size window of previously generated tokens.

\section{A Short Note}
\label{sec:shortnote}

This section is deliberately tiny.

\section{Cache Layout}
\label{sec:cache}

A key-value cache stores, for every layer, one vector per position for keys and
one for values. The layout follows the sequence structure; see
Section~\ref{sec:background} for the generation model and
Section~\ref{sec:shortnote} for the shortest block in this document.

\section{Eviction}
\label{sec:eviction}

Removing a contiguous range of positions frees the cells that range occupied.
The freed cells are reused by the next decode step.

\begin{theorem}
\label{thm:eviction}
A contiguous position range owned by one sequence can be removed in a single
operation.
\end{theorem}

\begin{proof}
\label{proof:eviction}
Every cell records which sequences it belongs to and what position it holds, so
a removal is fully described by a sequence identifier and a position range.
\end{proof}

\section{Conclusion}
\label{sec:conclusion}

Structured documents admit block-level eviction, which is the property this
project measures.

\end{document}